% =====================================================================
\chapter{Variables aleatorias discretas}
\label{tema:5}
% =====================================================================

Una variable aleatoria asigna un número a cada resultado de un experimento; su distribución resume todo lo que interesa de ella. Este tema introduce la función de masa, la esperanza y la varianza, demuestra la linealidad de la esperanza y la explota mediante variables indicadoras, y estudia las distribuciones discretas básicas: Bernoulli, binomial, geométrica, Poisson y categórica. Se siguen los capítulos~3 y~4 de Blitzstein y Hwang~\cite{bh}, el capítulo~2 de Bertsekas y Tsitsiklis~\cite{bt} (clases~5 a~7 de MIT 6.041~\cite{mit6041f10}), el capítulo~4 de Ross~\cite{ross} y el capítulo~18, «Random Variables», de~\cite{mcs}.

\section{Variables aleatorias y función de masa}

\begin{definicion}[Variable aleatoria discreta]
Una \emph{variable aleatoria} es una función $X\colon\Omega\to\R$. Es \emph{discreta} si toma un conjunto finito o numerable de valores. Su \emph{función de masa} es $p_X(x)=\PP(X=x)$, donde $\{X=x\}$ abrevia el suceso $\{\omega\in\Omega: X(\omega)=x\}$, y su \emph{función de distribución} es $F_X(x)=\PP(X\le x)$ para $x\in\R$.
\end{definicion}

La función de masa cumple $p_X(x)\ge0$ y $\sum_xp_X(x)=1$, porque los sucesos $\{X=x\}$ forman una partición de $\Omega$. Para todo conjunto $S\subseteq\R$, $\PP(X\in S)=\sum_{x\in S}p_X(x)$. La función de distribución es no decreciente, tiende a 0 en $-\infty$ y a 1 en $+\infty$, es continua por la derecha y, en el caso discreto, es escalonada, con saltos de altura $p_X(x)$ en cada valor $x$.

\addtocounter{definicion}{1}%

\section{Esperanza y linealidad}

\begin{definicion}[Esperanza]
La \emph{esperanza} (o media) de una variable discreta $X$ es
\[
\E[X]=\sum_{x}x\,p_X(x),
\]
siempre que la serie converja absolutamente. Es el centro de masas de la distribución.
\end{definicion}

\begin{proposicion}[Esperanza de una función: LOTUS]\label{prop:t5-lotus}
Si $X$ es discreta y $g\colon\R\to\R$, entonces $\E[g(X)]=\sum_xg(x)\,p_X(x)$. Si $\Omega$ es discreto, además $\E[X]=\sum_{\omega\in\Omega}X(\omega)\,\PP(\{\omega\})$.
\end{proposicion}

\begin{teorema}[Linealidad de la esperanza]\label{teo:t5-linealidad}
Si $X$ e $Y$ son variables aleatorias discretas definidas sobre el mismo espacio y con esperanza finita, y $a,b\in\R$, entonces
\[
\E[aX+bY]=a\,\E[X]+b\,\E[Y].
\]
No se requiere ninguna hipótesis de independencia.
\end{teorema}

\begin{definicion}[Variable indicadora]
La \emph{indicadora} de un suceso $A$ es la variable $I_A$ que vale 1 si ocurre $A$ y 0 si no.
\end{definicion}

\begin{proposicion}[Puente fundamental]\label{prop:t5-puente}
$\E[I_A]=\PP(A)$. En consecuencia, si $X=I_{A_1}+\dots+I_{A_n}$ cuenta cuántos de los sucesos $A_1,\dots,A_n$ ocurren, $\E[X]=\PP(A_1)+\dots+\PP(A_n)$.
\end{proposicion}

\addtocounter{definicion}{1}%

\section{Varianza e independencia}

\begin{definicion}[Varianza]
La \emph{varianza} de $X$ es $\Var(X)=\E\bigl[(X-\mu)^{2}\bigr]$, con $\mu=\E[X]$, y su \emph{desviación típica} es $\sigma_X=\sqrt{\Var(X)}$. Mide la dispersión de $X$ alrededor de su media, en unidades al cuadrado.
\end{definicion}

\begin{proposicion}\label{prop:t5-var}
$\Var(X)=\E[X^{2}]-\E[X]^{2}$ y, para $a,b\in\R$, $\Var(aX+b)=a^{2}\Var(X)$. En particular $\Var(X)\ge0$ implica $\E[X^{2}]\ge\E[X]^{2}$.
\end{proposicion}

\begin{definicion}[Independencia de variables]
Las variables discretas $X_1,\dots,X_n$ son \emph{independientes} si $\PP(X_1=x_1,\dots,X_n=x_n)=\prod_i\PP(X_i=x_i)$ para todos los valores $x_1,\dots,x_n$.
\end{definicion}

\begin{proposicion}\label{prop:t5-producto}
Si $X$ e $Y$ son discretas, independientes y con esperanza finita, $\E[XY]=\E[X]\E[Y]$. En consecuencia, si $X_1,\dots,X_n$ son independientes (basta dos a dos), $\Var(X_1+\dots+X_n)=\Var(X_1)+\dots+\Var(X_n)$.
\end{proposicion}

\section{Distribuciones discretas básicas}

\begin{definicion}[Bernoulli y binomial]
$X\sim\Ber(p)$, con $p\in[0,1]$, si $\PP(X=1)=p$ y $\PP(X=0)=1-p$: es la indicadora de un «éxito» de probabilidad $p$. $X\sim\Bin(n,p)$ si $X$ es el número de éxitos en $n$ ensayos de Bernoulli independientes con la misma probabilidad $p$.
\end{definicion}

\begin{teorema}\label{teo:t5-binomial}
Si $X\sim\Bin(n,p)$, entonces $\PP(X=k)=\binom nk p^{k}(1-p)^{n-k}$ para $k=0,\dots,n$, $\E[X]=np$ y $\Var(X)=np(1-p)$.
\end{teorema}

\begin{definicion}[Geométrica]
$X\sim\Geom(p)$, con $0<p\le1$, si $X$ es el número de ensayos de Bernoulli independientes de parámetro $p$ necesarios hasta obtener el primer éxito (incluido este). Su función de masa es $\PP(X=k)=(1-p)^{k-1}p$ para $k=1,2,\dots$
\end{definicion}

\begin{proposicion}\label{prop:t5-geom}
Si $X\sim\Geom(p)$: (i) $\PP(X>k)=(1-p)^{k}$; (ii) $\E[X]=1/p$ y $\Var(X)=(1-p)/p^{2}$; (iii) \emph{falta de memoria}: $\PP(X>m+k\mid X>m)=\PP(X>k)$ para enteros $m,k\ge0$.
\end{proposicion}

\begin{definicion}[Poisson]
$X\sim\Poi(\lambda)$, con $\lambda>0$, si $\PP(X=k)=e^{-\lambda}\lambda^{k}/k!$ para $k=0,1,2,\dots$ Se cumple $\E[X]=\Var(X)=\lambda$ (\cref{eje:5-poisson}).
\end{definicion}

La distribución de Poisson modela el número de sucesos raros entre muchas oportunidades independientes: erratas en un libro, llegadas de peticiones a un servidor en un minuto, apariciones de una palabra poco frecuente en un texto largo. El resultado siguiente lo justifica.

\begin{teorema}[Aproximación de Poisson a la binomial]\label{teo:t5-poisson}
Sea $\lambda>0$ y $p_n=\lambda/n$. Para cada $k\ge0$ fijo,
\(
\lim_{n\to\infty}\binom nk p_n^{k}(1-p_n)^{n-k}=e^{-\lambda}\frac{\lambda^{k}}{k!}.
\)
\end{teorema}

\begin{definicion}[Categórica]
Sea $\V=\{v_1,\dots,v_K\}$ un conjunto finito y $\bm p=(p_1,\dots,p_K)$ con $p_i\ge0$ y $\sum_ip_i=1$. Una variable $X$ con valores en $\V$ tiene distribución \emph{categórica}, $X\sim\Cat(\bm p)$, si $\PP(X=v_i)=p_i$. Cuando los valores no son números, en lugar de la esperanza de $X$ se usan las de sus indicadoras: $\E[I_{\{X=v_i\}}]=p_i$.
\end{definicion}

\begin{table}[h]
\centering
\begin{tabular}{@{}lllll@{}}
\toprule
Distribución & Valores & $\PP(X=k)$ & $\E[X]$ & $\Var(X)$ \\
\midrule
$\Ber(p)$ & $\{0,1\}$ & $p^{k}(1-p)^{1-k}$ & $p$ & $p(1-p)$ \\
$\Bin(n,p)$ & $\{0,\dots,n\}$ & $\binom nk p^{k}(1-p)^{n-k}$ & $np$ & $np(1-p)$ \\
$\Geom(p)$ & $\{1,2,\dots\}$ & $(1-p)^{k-1}p$ & $1/p$ & $(1-p)/p^{2}$ \\
$\Poi(\lambda)$ & $\{0,1,\dots\}$ & $e^{-\lambda}\lambda^{k}/k!$ & $\lambda$ & $\lambda$ \\
Uniforme en $[n]$ & $\{1,\dots,n\}$ & $1/n$ & $(n+1)/2$ & $(n^{2}-1)/12$ \\
\bottomrule
\end{tabular}
\caption{Distribuciones discretas básicas.}
\label{tab:t5-discretas}
\end{table}

\addtocounter{definicion}{1}%

\addtocounter{definicion}{1}%

